\subsection{Alcance e desvio dos trajetos}\label{sec:F-EFI-2}
Em $G_E$, o Dijkstra limita a busca a $T$ segundos e exibe as paradas
fechadas, com custo $O((V'+E')\log(E'+2))$ na parte visitada. No pico da
manhã, o alcance a partir da parada 14788851 cresceu conforme o quadro abaixo.
\begin{center}\small
\begin{tabular}{lrrrr}
Tempo disponível (min) & 15 & 30 & 45 & 60\\
Paradas alcançadas & 68 & 603 & 1.911 & 4.120
\end{tabular}
\end{center}
Em 12 pares sorteados (semente 2026), o maior desvio foi 2,22: para a viagem
14793286 a 14785528, o trajeto mais rápido percorreu 1,47 km para vencer
0,67 km em linha reta. O comprimento soma as distâncias entre paradas
consecutivas e pode subestimar curvas das ruas.
